\documentclass[11pt]{article}
\usepackage[utf8]{inputenc}
\usepackage[T1]{fontenc}
\usepackage[italian]{babel}
\usepackage{geometry}
\geometry{a4paper, margin=1in}
\usepackage{xcolor}
\usepackage{listings}
\usepackage{hyperref}

\usepackage{tcolorbox}
\tcbuselibrary{breakable}
\begin{document}

\begin{tcolorbox}[colback=blue!5!white,colframe=blue!75!black,title=\textbf{DOMANDA}, breakable]
Sommatrice UDP e perdita di pacchetti\\
9) Usare la sommatrice su due macchine distinte provando, sulla macchina del client, a staccare il cavo di rete prima di un invio di un dato, ad es.
\begin{itemize}
    \item Digitare "2345" + INVIO
    \item Digitare "5187" + INVIO
    \item Staccare il cavo
    \item Digitare "2" + INVIO
    \item Riattaccare il cavo e aspettare 30 sec che il sistema operativo si riassesti
    \item Digitare "1" + INVIO
    \item "0"
\end{itemize}
\begin{itemize}
    \item Che somma leggo? E' corretta?
    \item Se il server gira su una macchina Linux (anche una macchina virtuale) e comunica con il client attraverso una rete WiFi oppure in loopback è possibile simulare la perdita di pacchetti bloccando e, dopo poco tempo, sbloccando la ricezione di pacchetti con i seguenti comandi:
    \begin{itemize}
        \item PER BLOCCARE: \texttt{sudo iptables -A INPUT -p udp --dport <porta\_del\_server> -j DROP}
        \item PER SBLOCCARE: \texttt{sudo iptables -D INPUT -p udp --dport <porta\_del\_server> -j DROP}
        \item Si consiglia di prepararsi questi comandi in una shell per poterli dare velocemente
    \end{itemize}
\end{itemize}
\end{tcolorbox}

\begin{tcolorbox}[colback=green!5!white,colframe=green!75!black,title=\textbf{RISPOSTA}, breakable]
Questo esercizio dimostra in modo pratico una delle caratteristiche fondamentali del protocollo \texttt{UDP}: la sua inaffidabilità (best-effort delivery) e l'assenza di connessione.

Rispondendo alle domande:

\begin{itemize}
    \item \textbf{Che somma leggo?} Leggerai la somma 7533 (derivata da 2345 + 5187 + 1).
    \item \textbf{È corretta?} Dal punto di vista dell'utente non è corretta, in quanto manca il 2 (la somma attesa sarebbe dovuta essere 7535). Dal punto di vista del protocollo \texttt{UDP}, invece, è un comportamento del tutto normale.
\end{itemize}

Ecco cosa succede nel dettaglio e il perché di questo risultato:

\begin{enumerate}
    \item I primi due pacchetti (2345 e 5187) vengono inviati e raggiungono il server, che inizia ad accumulare la somma.
    \item Quando scolleghi il cavo (o attivi la regola di iptables per far scartare (DROP) i pacchetti) e digiti 2, la funzione \texttt{UDPSend} nel codice del client non restituisce alcun errore. Per il client l'invio ha avuto successo, perché \texttt{UDP} non prevede alcun meccanismo di conferma di ricezione (ACK). Il pacchetto contenente il "2" viene semplicemente scartato dalla rete e perso per sempre (packet loss).
    \item Quando riattacchi il cavo e invii 1 e 0, la rete è di nuovo operativa. Questi pacchetti raggiungono il server che li somma ai precedenti.
    \item Quando arriva lo 0, il server risponde inviando il totale accumulato (7533) al client.
\end{enumerate}

Questo esercizio mette in luce che, qualora si utilizzi \texttt{UDP} per applicazioni che necessitano di garanzie di consegna, è necessario implementare i controlli di ricezione, numerazione dei pacchetti e ritrasmissione direttamente a livello applicativo, poiché \texttt{UDP} a livello di trasporto non offre nulla di tutto ciò.
\end{tcolorbox}

\end{document}
